\chapter{Agent 产品与交互设计}
\label{chap:product}

\begin{takeaway}[学习目标]
把 Agent 能力转成可理解、可控制的用户体验；设计委托、进度、确认、纠错和信任机制；用业务结果与人工成本衡量产品价值。
\end{takeaway}

\section{用户购买的是结果，不是自主性}

“全自动 Agent”不是价值主张。用户关心的是能否更快、更准、更省心地完成任务，以及出错时能否发现和修正。产品定义应写成“为谁，在什么情境，把哪项指标从多少改善到多少”。

\section{把任务边界展示清楚}

好的委托界面让用户在开始前知道：目标、输入来源、输出、预计时间和费用、可访问范围、会产生的外部动作、需要确认的节点。开放聊天框适合探索，但高频业务应逐渐沉淀为表单、模板和可复用任务。

系统对含糊目标可以先生成“任务理解卡”：范围、假设、计划和完成条件。用户修改结构化字段比反复聊天纠正整段 prompt 更有效。

\section{渐进式自主}

同一产品可提供四档控制：

\begin{enumerate}
  \item \textbf{建议}：只给方案，不执行。
  \item \textbf{草稿}：执行可逆步骤，外部动作保留为草稿。
  \item \textbf{确认后执行}：在关键边界请求批准。
  \item \textbf{受限自动}：对已授权的低风险重复任务自动运行。
\end{enumerate}

自主级别应按任务和工具决定，而不是一个全局开关。随着历史成功和用户授权可以提高，但始终提供暂停、撤销和查看活动的入口。

\section{进度应该反映真实状态}

不要展示虚假的“思考中”。用户需要看到可验证的阶段：正在核验 3/5 个来源、等待网站响应、已生成草稿、需要批准发送。长任务给预计范围、已消耗预算和下一检查点。

内部 chain-of-thought 不适合作为产品解释。更有用的是行动摘要：使用了哪些来源、调用了哪些工具、为什么需要确认、哪些信息仍不确定。

\section{失败体验决定信任}

失败界面回答四件事：发生了什么；哪些步骤已经完成；是否产生副作用；用户现在能做什么。提供从检查点继续、修改输入、换数据源、下载部分结果和转人工的操作。

“出错了，请重试”会抹掉 Agent 长任务的上下文价值。错误文案应关联任务和错误 ID，但不要泄露内部密钥、堆栈或其他租户信息。

\section{确认与撤销的细节}

确认框展示动作对象、差异、接收者、金额、可见范围和不可逆后果。批量操作提供筛选与逐项取消。确认后资源若变化，应使旧批准失效并重新展示 diff。

撤销应靠产品能力实现，而非承诺“AI 会处理”。对文件使用版本历史，对草稿支持删除，对业务记录使用补偿动作。不可撤销动作要更早预览、更晚提交。

\section{信任来自证据和校准}

展示来源、时间、测试结果和实际工具回执。系统应该在证据不足时降低确定性，而不是总用自信语气。置信度最好与可操作行为绑定：低置信度请求澄清或人工复核，高置信度仍受权限规则限制。

用户反馈不能只有点赞。提供“事实错误、漏项、未遵循要求、工具选错、太慢、太贵、越权”等可诊断类别，并关联 trace。

\section{产品指标树}

\begin{table}[htbp]
\centering
\begin{tabularx}{\textwidth}{>{\bfseries}p{0.2\textwidth} X X}
\toprule
层级 & 核心问题 & 指标示例 \\
\midrule
业务结果 & 任务是否创造价值？ & 处理时长、转化、节省人时 \\
任务质量 & 结果是否可直接使用？ & 完成率、一次通过率、修订率 \\
控制与信任 & 用户是否敢委托？ & 批准率、撤销率、越权事件 \\
效率 & 系统代价是否合理？ & 延迟、成本、人工接管分钟数 \\
留存 & 是否形成稳定工作流？ & 周活任务、模板复用、留存 \\
\bottomrule
\end{tabularx}
\end{table}

自动化率不能单独作为北极星。把需要人工修正的工作藏到流程末端，可能让自动化率上升、总耗时反而增加。

\section{从辅助驾驶到自动化}

新场景先运行 shadow mode：Agent 产生建议但不执行，与人工结果对比。达到质量门槛后开放草稿；再开放确认执行；只有低风险、高频、稳定任务进入受限自动化。每一级都需要退出和回滚条件。

\begin{practice}[画一个关键流程]
为你的 Agent 画出从委托、运行、确认、完成到失败恢复的界面流程。为每个页面写出用户能看到的真实状态和可执行动作。找五位目标用户完成任务，记录他们在哪一步不敢继续。
\end{practice}

\begin{checkpoint}
\begin{itemize}
  \item \checkbox 产品价值用业务结果描述，而不是“更自主”。
  \item \checkbox 用户能看到范围、权限、进度、证据和费用。
  \item \checkbox 失败后能看到已完成部分、副作用与恢复选项。
  \item \checkbox 自主级别按风险逐步开放，并有回滚门槛。
\end{itemize}
\end{checkpoint}

